\begin{luacode*}

function gerar_estatisticas_csv(input_csv, output_csv, totais_csv)

-- converter data + hora
function parse_datetime(data,hora)
  if not data or data=="" then return nil end
  if not hora or hora=="" then hora="00:00:00" end

  local d,m,y = data:match("(%d%d)/(%d%d)/(%d%d%d%d)")
  local H,M,S = hora:match("(%d%d):(%d%d):?(%d*)")

  if not d then return nil end

  return os.time{
    year=tonumber(y),
    month=tonumber(m),
    day=tonumber(d),
    hour=tonumber(H) or 0,
    min=tonumber(M) or 0,
    sec=tonumber(S) or 0
  }
end


-- parser CSV
function split_csv(line)

  local res,pos={},1

  while pos<=#line do
    local c=line:sub(pos,pos)

    if c=='"' then
      local first=pos+1
      local last=line:find('"',first)

      while line:sub(last+1,last+1)=='"' do
        last=line:find('"',last+2)
      end

      table.insert(res,line:sub(first,last-1))
      pos=last+2

    else
      local next=line:find(',',pos) or (#line+1)
      table.insert(res,line:sub(pos,next-1))
      pos=next+1
    end
  end

  for i=1,#res do
    res[i]=res[i]:match("^%s*(.-)%s*$")
  end

  return res
end


-- registrar guia
function registrar_guia(nome,total_part,duracao,guias)

  if not nome or nome=="" then return end

  nome=nome:match("^%s*(.-)%s*$")

  if not guias[nome] then
    guias[nome]={atividades=0,participantes=0,horas=0}
  end

  guias[nome].atividades=guias[nome].atividades+1
  guias[nome].participantes=guias[nome].participantes+total_part
  guias[nome].horas=guias[nome].horas+duracao

end


local file=io.open(input_csv,"r")
if not file then
  texio.write_nl("Erro ao abrir "..input_csv)
  return
end

local guias={}

local total_atividades=0
local total_participantes=0
local total_horas=0

file:read()

for line in file:lines() do

  local col=split_csv(line)

  local guia1=col[2]
  local guia2=col[15]

  local participantes=col[10]

  local data_ini=col[6]
  local data_fim=col[7]

  local hora_ini=col[8]
  local hora_fim=col[14]

  -- contar participantes
  local total_part=0

  if participantes and participantes~="" then
    for nome in participantes:gmatch("[^,]+") do
      nome=nome:match("^%s*(.-)%s*$")
      if nome~="" then total_part=total_part+1 end
    end
  end

  -- calcular duração
  local inicio=parse_datetime(data_ini,hora_ini)

  local fim
  if data_fim and data_fim~="" then
    fim=parse_datetime(data_fim,hora_fim)
  else
    fim=parse_datetime(data_ini,hora_fim)
  end

  local duracao=0
  if inicio and fim then
    duracao=(fim-inicio)/3600
  end

  -- totais gerais
  total_atividades=total_atividades+1
  total_participantes=total_participantes+total_part
  total_horas=total_horas+duracao

  -- registrar guias
  registrar_guia(guia1,total_part,duracao,guias)

  if guia2 and guia2~=guia1 then
    registrar_guia(guia2,total_part,duracao,guias)
  end

end

file:close()


-- ranking
local ranking={}

for g,d in pairs(guias) do
  table.insert(ranking,{
    guia=g,
    atividades=d.atividades,
    participantes=d.participantes,
    horas=d.horas
  })
end

table.sort(ranking,function(a,b)
  if a.atividades~=b.atividades then
    return a.atividades>b.atividades
  elseif a.participantes~=b.participantes then
    return a.participantes>b.participantes
  else
    return a.horas>b.horas
  end
end)


-- escrever CSV ranking
local out=io.open(output_csv,"w")
out:write("Guia,Atividades,Participantes,Horas\n")

for _,item in ipairs(ranking) do
  out:write(string.format(
  "%s,%d,%d,%.2f\n",
  item.guia,
  item.atividades,
  item.participantes,
  item.horas))
end

out:close()


-- escrever totais
if totais_csv then
  local out2=io.open(totais_csv,"w")
  out2:write("Atividades,Participantes,Horas\n")
  out2:write(string.format(
  "%d,%d,%.2f\n",
  total_atividades,
  total_participantes,
  total_horas))
  out2:close()
end

end


-- texto automático
function texto_estatisticas_totais(csv)

  local f=io.open(csv,"r")
  if not f then
    tex.sprint("Dados estatísticos indisponíveis.")
    return
  end

  f:read()
  local line=f:read()
  f:close()

  local a,p,h=line:match("([^,]+),([^,]+),([^,]+)")

  tex.sprint(string.format(
  "Durante o período coberto por este relatório, foram realizadas \\textbf{%s} pranchetas oficiais, mobilizando \\textbf{%s} participantes e totalizando \\textbf{%.0f} horas de atividades de montanha.",
  a,p,tonumber(h)))

end

\end{luacode*}
